\section{Sixth order: a shape that cannot be removed}
\label{sec:sixth}

The fourth-order shape data closes in a five-element span. The question
this section answers is whether the sixth order stays inside it. It does
not, and the proof is a divisibility argument that needs no knowledge of
the sixth-order dynamics at all.

\subsection{Conventions}

In the Bloch variables $z_j$ write
\begin{equation}
  a_j = 2 - z_j - z_j^{-1},
  \qquad q = \sum_j a_j,
  \qquad e_2 = \sum_{i<j} a_i a_j,
  \qquad e_3 = a_1 a_2 a_3 ,
\end{equation}
and let $U = \psi\psi^{*}$ with $\psi^{*}\psi = q$, $P_c = U/q$,
$Q_c = I - P_c$, $S = L_{\downarrow} - 4I$, and $R$ the cross-plane part
of $S$.

Two warnings, both of which have caused errors here. The invariants
$e_2, e_3$ are Bloch invariants, \emph{not} representation Casimirs;
and $U$ is the up-Laplacian, not its normalized projector, so its kernel
is the complementary plane rather than the carrier line. Cleared
identities hold at $q = 0$; normalized carrier formulas require
$q > 0$.

\subsection{The complete folded operator}

Let $P_E$ be the isolated electric model-space projector,
$D = Q_E/(E_0 - H_0)$, and $V_0$ the perturbation. If
$P_E V_0 P_E = a P_E$, set $V = V_0 - aI$ \emph{on the whole Hilbert
space}: shifting the complementary block as well is essential. With
$\chi_0 = P_E$, intermediate normalization gives the exact recursions
\begin{equation}
  K_n = P_E V \chi_{n-1},
  \qquad
  \chi_n = D V \chi_{n-1} - D \sum_{j=1}^{n-1} \chi_{n-j} K_j ,
\end{equation}
and with $M = P_E + \chi^{*}\chi$ the canonical Hermitian effective
operator is $M^{1/2} K M^{-1/2}$. Writing
\begin{equation}
  A_m = P_E V (DV)^{m-1} P_E,
  \qquad
  C_2 = P_E V D^3 V P_E,
\end{equation}
and $B_m$ for the sum of $A_m$ with exactly one of its $m-1$ slots
replaced by $D^2$:

\begin{theorem}[The eighteen-word sixth-order fold]
\label{thm:folds}
\begin{equation}
  H_6 = A_6 - \tfrac12\{A_4,B_2\} - \tfrac12\{A_3,B_3\}
        - \tfrac12\{A_2,B_4\} + \tfrac12\{A_2^2,C_2\}
        + \tfrac38\{A_2,B_2^2\} + \tfrac14 B_2 A_2 B_2 .
  \label{eq:F6}
\end{equation}
Expanded, this is exactly eighteen words in $P_E, V, D$. The same
recursion reproduces the established $H_4 = A_4 - \tfrac12\{A_2,B_2\}$.
\end{theorem}
\gid{RESULT:G9\_HERMITIAN\_FOLDED\_SUPPORT}

The $\{A_3,B_3\}$ term is present and it matters: dropping odd electric
histories gives an incomplete formula. No odd order is dropped anywhere
in the recursion. Independent rational matrix checks verify both the
invariant-subspace equation and the normalization through degree six,
with a rank-two $P_E$ and non-commuting effective coefficients.

\notasserted{Equation~\eqref{eq:F6} is the universal folded operator
support in this convention. It does not evaluate Wilson plaquette matrix
elements, establish a specific support, or carry out connected
multi-plaquette rooted subtraction. Those require the actual electric
and Haar implementation on each support.}

\subsection{Exact word symbols}

In a polynomial carrier gauge, exact symmetrization gives each word its
cleared symbol in $\Q[q, e_2, e_3]$. The table below is computed, not
inferred, and all forty words of length zero through three agree
coefficient by coefficient with independently implemented spatial
Laurent operators.

\begin{center}
\begin{tabular}{@{}ll@{\qquad}ll@{}}
\toprule
Word & $\sigma(W)$ & Word & $\sigma(W)$ \\
\midrule
$I$        & $q$          & $RR$        & $q e_2 + 3 e_3$ \\
$U$        & $q^2$        & $RUR$       & $4 e_2^2$ \\
$S$        & $-4q$        & $RSR$       & $(q-4)(q e_2 + 3e_3) - 4e_2^2$ \\
$S^2$      & $16q$        & $RRR$       & $-2e_2^2 - 2q e_3$ \\
$R$        & $-2e_2$      & $URR,\,RRU$ & $q^2 e_2 + 3 q e_3$ \\
$UR,\,RU$  & $-2q e_2$    &             & \\
\bottomrule
\end{tabular}
\end{center}

A word census must retain coefficients and combine them before testing
shape membership: $\sigma(RUR + 2\,RRR) = -4q e_3$, although both
individual words contain $e_2^2$. Support is not unique modulo operator
identities, and the invariant object being tested is the complete
carrier polynomial. Missing direct input is represented as unknown, not
as zero.

\subsection{The band coefficient at sixth order}

At fixed nonzero momentum, with scalar anchors removed,
\begin{equation}
  H_{\mathrm{eff}}(u) = u^2 t_3\, q\, Q_c + u^3 H_3 + u^4 H_4
    + u^5 H_5 + u^6 H_6 + O(u^7).
\end{equation}
The recorded $\SU(3)$ third-order factorization gives
$Q_c H_3 P_c = 0$, so the off-diagonal block first appears at order
$u^2$ in the scaled operator, from $H_4$. Expanding the Schur complement
for the scalar carrier branch,
\begin{equation}
  \lambda_6 = \frac{\sigma(H_6)}{q}
    - \frac{q\,\sigma(H_4^2) - \sigma(H_4)^2}{t_3\, q^3}.
  \label{eq:S6}
\end{equation}
This is not an assumption that $H_5$ vanishes. $H_5$ first couples to
$H_4$ at order $u^7$ in $H_{\mathrm{eff}}$, as does the $uH_3$
correction to the complementary inverse, and unknown off-diagonal $H_6$
contributes later still. The decoupling of $H_3$ is an explicit recorded
input rather than an assumption, and any scalar vacuum subtraction is
part of $H_6$ and changes nothing.

Evaluating all twenty-five ordered pairs from the actual $H_4$ support
$\{I,U,S,S^2,R\}$, including their projected subtraction, gives zero for
twenty-four of them. Only $RR$ survives, in agreement with direct
composition of the full assembled spatial kernel. With
\begin{equation}
  \kappa = \frac{4C_{\mathrm{shp}}^2}{t_3}
  = \frac{169924002729806205028788409}{619343593697933825385600000} > 0,
  \qquad t_3 = \tfrac{5}{612},
\end{equation}
the known dynamical mixing support is
$-(\kappa/q)\,RR + (\kappa/q^2)\,RUR$, with normalized expectation
\begin{equation}
  \lambda_6^{\mathrm{mix}}
  = -\kappa\left(\frac{e_2}{q} + \frac{3e_3}{q^2}
      - \frac{4e_2^2}{q^3}\right).
\end{equation}

\notasserted{The $RUR$ here arises from inserting $P_c = U/q$ in the
projected subtraction. It is not evidence that a direct six-insertion
electric history carries the local word $RUR$, and the ratio $1:3:-4$
belongs to this mixing term rather than to the full coefficient.}

\subsection{Non-cancellation}

\begin{theorem}[An extra rational shape survives every local direct term]
\label{thm:noncancellation}
Assume the order-six electric-shell effective operator is a
finite-range, translation-covariant stencil with Laurent-polynomial
entries in this basis, and write $F = \sigma(H_6)$. Then \eqref{eq:S6}
gives
\begin{equation}
  q^3 \lambda_6 = q^2 F - \kappa\big(q^2 e_2 + 3q e_3 - 4e_2^2\big),
  \label{eq:N6}
\end{equation}
and $q$ does not divide the right-hand side, for \emph{any} Laurent
polynomial $F$.
\end{theorem}

\begin{proof}
Modulo $q$, the right-hand side of \eqref{eq:N6} is $4\kappa e_2^2$.
Evaluate at
\begin{equation}
  z_1 = z_2 = -1, \qquad z_3 = 5 + 2\sqrt{6},
\end{equation}
all nonzero, so the point lies in the Laurent ring. Since
$z_3^{-1} = 5 - 2\sqrt6$, we get $a = (4, 4, -8)$, hence $q = 0$ and
$e_2 = -48$. The value there is $4\kappa \cdot 48^2 = 9216\,\kappa \neq 0$,
and a Laurent polynomial divisible by $q$ would vanish at that point.
\end{proof}
\gid{RESULT:G9\_LOCAL\_SIXTH\_NONCANCELLATION}

\noindent\Tone{}. Consequently $\lambda_6$ cannot lie in the
fourth-order shape span $\{1, q, e_2, 4e_2/q, e_3/q\}$, and no local
direct $RUR$ term removes the extra rational shape: a direct $h\,RUR$
contributes $4h e_2^2/q$, which cannot cancel an $e_2^2/q^3$ component.
For symmetric $F$ the remainder modulo $q^2$ is exactly
$4\kappa e_2^2 - 3\kappa q e_3$, so both $e_3/q^2$ and $e_2^2/q^3$ are
present. The $e_2/q$ term alone \emph{can} be cancelled, by
$F = \kappa e_2$ --- for instance a local $H_6 = -(\kappa/2)R$ --- so
the $1:3:-4$ ratio is not a constraint on all sixth-order shapes.

Note what the proof does not use: any knowledge of $F$. The extra shape
is forced by the fourth-order data alone, through the Schur complement,
whatever the sixth-order dynamics turns out to be.

\subsection{A withdrawn derivation}

An earlier version of this work enumerated closed rooted walks --- $90$
four-step unit-edge walks, of which $66$ retrace immediately and $24$
are planar squares, together with $192$ non-retracing six-step walks
using all three axes --- and read a sixth-order amplitude off that
census.

The derivation is withdrawn. Edge walks are not ordered plaquette
insertions with electric states, Haar weights, projectors and folds, and
no map between them was supplied; the census therefore determines no net
coefficient of any of the candidate words. The old amplitude formula was
assigned rather than derived. The old carrier function used
$\sigma(R) = 0$ and $\sigma(U) = q$, where the actual unnormalized
symbols are $-2e_2$ and $q^2$; it returned a nonzero expression for a
defect it separately reported as zero. The implementation now computes
the symbols and tests the defect, and the old interface raises an
explicit missing-dynamics error rather than returning a number.

What is withdrawn is a claimed derivation, not the possibility of a
physical amplitude of that shape. Theorem~\ref{thm:noncancellation} is
independent of the census entirely. Computing the direct term on a
connected multi-plane support is Problem~\ref{prob:g9}, and it is the
cheapest decisive computation currently open in this program.
